\section{Visão Geral da Organização}

\subsection{Estrutura de alto nível}

O repositório está organizado por responsabilidades técnicas. A raiz contém a solução .NET, a configuração global, a infraestrutura local, a interface web, os dados de suporte, os scripts, os testes e a documentação.

Uma vista simplificada da raiz é a seguinte:

\begin{verbatim}
NatureProtector/
|-- data/
|-- docs/
|-- infra/
|-- scripts/
|-- src/
|-- tests/
|-- webUI/
|-- docker-compose.yml
|-- NatureProtector.sln
|-- NuGet.Config
|-- coverage.runsettings
`-- README.md
\end{verbatim}

Esta organização separa o código executável dos testes, dos dados, dos scripts de automação, da documentação e dos artefactos de infraestrutura.

\subsection{Principais blocos do repositório}

\begin{longtable}{L{0.28\textwidth} L{0.62\textwidth}}
\toprule
\textbf{Pasta/Ficheiro} & \textbf{Função principal} \\
\midrule
\endfirsthead
\toprule
\textbf{Pasta/Ficheiro} & \textbf{Função principal} \\
\midrule
\endhead
\repoPath{src/} & Código fonte dos projetos .NET, incluindo domínio, hosts, API, persistência e infraestrutura. \\
\repoPath{tests/} & Testes unitários, testes de integração e testes por módulo. \\
\repoPath{webUI/} & Interface web usada para monitorização, execução de cenários, diagnósticos e evidência. \\
\repoPath{docs/} & Documentação técnica, arquitetura, contratos, evidência, relatório, poster e materiais auxiliares. \\
\repoPath{infra/} & Configuração da infraestrutura local, incluindo PostgreSQL, Grafana e scripts de infraestrutura. \\
\repoPath{scripts/} & Scripts de setup, execução, cenários, evidência, documentação, testes e apoio ao desenvolvimento. \\
\repoPath{data/} & Dados de suporte, baseline, manifestos de datasets e manifestos de cenários. \\
\repoPath{docker-compose.yml} & Definição dos serviços locais usados pela baseline. \\
\repoPath{NatureProtector.sln} & Solução .NET principal. \\
\repoPath{coverage.runsettings} & Configuração usada para recolha de cobertura de testes. \\
\bottomrule
\end{longtable}

\subsection{Separação entre código, documentação e evidência}

O repositório não contém apenas código fonte. Contém também documentos de arquitetura, mapas de implementação, contratos, dados, scripts de execução e evidência runtime. Esta separação permite consultar não só a implementação, mas também a forma como a implementação foi validada.

A documentação encontra-se sobretudo em \repoPath{docs/}. A evidência de execução e validação é guardada em \repoPath{docs/evidence/}, enquanto scripts reprodutíveis de apoio se encontram em \repoPath{scripts/}. Esta organização facilita a auditoria do projeto e reduz a dependência de explicações orais.

\subsection{Artefactos gerados}

Algumas pastas podem existir localmente por serem geradas durante build, testes, coverage, execução da UI ou execução de scripts. Exemplos comuns são:

\begin{itemize}
    \item \repoPath{bin/} e \repoPath{obj/};
    \item \repoPath{.testbin/};
    \item \repoPath{TestResults/};
    \item \repoPath{coveragereport_core/};
    \item \repoPath{node_modules/};
    \item \repoPath{__pycache__/};
    \item ficheiros auxiliares gerados por LaTeX, como \repoPath{.aux}, \repoPath{.log}, \repoPath{.fls} ou \repoPath{.toc}.
\end{itemize}

Estes artefactos não devem ser confundidos com a estrutura fonte do repositório e devem estar cobertos por \repoPath{.gitignore}, salvo quando a equipa decidir versionar explicitamente algum output final.
